\section{Introduction}
\label{sec:introduction}

Yayoi Kusama (born 1929) is one of the most influential living artists,
renowned for her obsessive deployment of polka dots across paintings,
sculptures, and immersive installations. Her visual language is built
upon four signature traits: geometrically precise hard-edged dots,
multi-scale coexistence of dot sizes ranging from sub-millimeter specks
to large medallions, ultra-pure high-chroma color palettes applied in
flat fields, and \emph{horror vacui}---a compulsion to cover every
available surface with dots, leaving no void untouched
\citep{moma_kusama,archdaily_kusama}. The polka dot, in
Kusama's universe, is simultaneously a decorative motif, a signifier of
infinity, and a mechanism for dissolving the boundary between the self
and the environment.

Translating this artistic language into a computational framework
presents a unique set of challenges. Classical stippling algorithms,
pioneered by \citet{secord2002weighted} and grounded in Centroidal
Voronoi Tessellation (CVT) and Lloyd relaxation \citep{lloyd1982least},
excel at producing blue-noise point distributions modulated by image
density. However, these methods map dot size and density to a single
scalar field---typically luminance: dark regions receive larger or
denser dots, bright regions receive fewer and smaller ones. This
one-dimensional mapping is fundamentally at odds with Kusama's
practice. A pumpkin painted by Kusama does not merely have dots of
varying size due to shading; the dots themselves carry a sense of
volume, of near and far, of the object inhabiting space. Capturing
this requires a control signal that luminance alone cannot provide.

Monocular depth estimation supplies exactly such a signal. The field
has matured rapidly, and recent unified real-time backbones---most
notably the YOLO26 family \citep{ultralytics2025yolo26}---now produce
dense metric depth maps at interactive rates on consumer GPUs. Three
considerations motivate our choice of YOLO26 over alternatives such as
MiDaS or DepthAnything: (i) a single checkpoint family also provides
semantic segmentation from the same feature backbone, guaranteeing
that the geometric and semantic control signals are spatially
consistent by construction; (ii) inference times of a few seconds at
$1024 \times 1024$ resolution fit an artist-in-the-loop workflow; and
(iii) the metric (rather than disparity) output simplifies the
normalization step that drives our radius mapping. We emphasize,
however, that the proposed pipeline is \emph{backbone-agnostic}: any
monocular depth estimator and any semantic segmentor can be
substituted without algorithmic changes.

In this paper, we propose a \emph{depth-aware polka dot art generation
pipeline} that integrates monocular depth estimation and semantic
segmentation with Weighted Voronoi Stippling. The central mathematical
component is a depth-to-radius mapping:

\begin{equation}
    r(x,y) = \Bigl[\, r_{\min} + (r_{\max} - r_{\min}) \cdot (1 - \hat{d})^\alpha \,\Bigr] \cdot (1 + \varepsilon),
    \label{eq:depth_radius}
\end{equation}

\noindent where $\hat{d}(x,y) \in [0,1]$ is the normalized depth (0 =
nearest, 1 = farthest), $\alpha$ is a curvature exponent controlling
the depth-size falloff, and $\varepsilon \sim \mathcal{N}(0,
\sigma^2)$ is a multiplicative Gaussian jitter that injects the
organic variation characteristic of hand-painted dots. Near
objects---which occupy more of the visual field and carry greater
perceptual weight---receive large, prominent dots, while distant
surfaces recede into fields of small, sparse dots. The result is a
polka dot rendering that conveys three-dimensional structure through
dot size alone, without resorting to perspective projection or
shading.

We further leverage semantic segmentation to enable \emph{region-aware
coloring}: foreground objects and background regions receive different
dot colors within the same duotone palette, creating a figure--ground
separation that echoes Kusama's practice of placing contrasting
colored dots on different surfaces within a single composition.

Two design decisions differentiate the pipeline from a naive
combination of off-the-shelf components. First, the radius range
demanded by the depth mapping (roughly 1--15\,px) is far wider than
what a single-layer stipple can support without blob merging. We
therefore decompose the rendering into a three-layer architecture:
(1) a fine stipple layer of small, densely packed dots whose density
follows a composite depth-luminance field; (2) a depth-emphasis layer
of large dots placed exclusively in near regions, pruned against the
fine layer; and (3) an edge and background layer providing structural
emphasis and horror-vacui coverage. Second, a dedicated cross-layer
overlap-pruning step enforces a minimum pairwise distance between the
large emphasis dots and the fine stipple, which we show
(Section~\ref{sec:experiments}) to be the single most important
component for preserving dot separation.

We summarize our contributions as follows:
\begin{enumerate}
    \item A depth-to-radius mapping function (Eq.~\ref{eq:depth_radius})
    that encodes scene depth into dot size with controlled
    multiplicative randomization, enabling three-dimensional perception
    in the inherently flat polka dot medium.
    \item A three-layer stippling architecture with cross-layer overlap
    pruning that maintains dot separation (maximum blob ratio below
    0.65\%) while supporting a wide dot radius range (0.5--15.4\,px).
    \item Semantic-aware dual-color rendering using a unified YOLO26
    backbone, assigning distinct dot colors to foreground and
    background regions within a Kusama-style palette.
    \item A \emph{learnable} depth-to-radius network trained via a
    self-supervised perceptual stippling loss, providing a data-driven
    alternative to the fixed power-law mapping that adapts across scene
    types without manual parameter tuning.
\end{enumerate}

The remainder of this paper is organized as follows.
Section~\ref{sec:related} reviews related work in stippling, depth
estimation, and neural style transfer. Section~\ref{sec:method} presents
the mathematical formulation and pipeline architecture in detail.
Section~\ref{sec:experiments} reports experimental results and quality
metrics. Section~\ref{sec:discussion} discusses artistic implications,
limitations, and future directions. Section~\ref{sec:conclusion}
concludes the paper. This work is part of a broader research program on
\emph{AI for Art Process}---using AI to understand and intervene in the
creative \emph{process} of art-making, complementary to work on
real-time painting coaching and painting-to-video animation.
